\documentclass[11pt]{amsart}

\usepackage[UTF8,scheme=plain,fontset=fandol]{ctex}
\usepackage{microtype}
\usepackage{amsmath,amssymb,mathtools}
\usepackage{booktabs}
\usepackage{tabularx}
\usepackage{longtable}
\usepackage{listings}
\usepackage{xcolor}
\usepackage{hyperref}

\lstset{
  basicstyle=\ttfamily\small,
  breaklines=true
}

\hypersetup{
  colorlinks=true,
  linkcolor=blue!55!black,
  citecolor=blue!55!black,
  urlcolor=blue!55!black,
  pdftitle={trace-73 Cappell--Shaneson skein-lasagna 阻碍（中文版）},
  pdfauthor={Anonymous repository audit}
}

\newtheorem{theorem}{定理}[section]
\newtheorem{proposition}[theorem]{命题}
\newtheorem{lemma}[theorem]{引理}
\newtheorem{corollary}[theorem]{推论}
\theoremstyle{definition}
\newtheorem{definition}[theorem]{定义}
\newtheorem{hypothesis}[theorem]{假设}
\theoremstyle{remark}
\newtheorem{remark}[theorem]{注}
\newtheorem{externalresult}[theorem]{外部结果}

\newcommand{\Q}{\mathbb{Q}}
\newcommand{\Z}{\mathbb{Z}}
\newcommand{\Sfour}{S^{4}}
\newcommand{\SN}{\mathcal{S}^{N}_{0}}
\newcommand{\Stwo}{\mathcal{S}^{2}_{0}}
\newcommand{\KhR}{\operatorname{KhR}}
\newcommand{\HH}{\operatorname{HH}_{0}}
\newcommand{\Span}{\operatorname{Span}}
\newcommand{\eps}{\varepsilon}
\newcommand{\Sh}{\operatorname{Sh}}
\newcommand{\im}{\operatorname{im}}
\newcommand{\Bur}{\mathrm{Bur}}
\newcommand{\MWW}{\mathrm{MWW}}
\newcommand{\sha}{\operatorname{SHA256}}
\newcommand{\Partial}{\textbf{部分}}
\newcommand{\Discharged}{\textbf{已证}}
\newcommand{\Unused}{\textbf{未用}}
\newcommand{\hash}[2]{\shortstack[l]{\texttt{#1}\\\texttt{#2}}}

\title[trace-73 阻碍]
{trace-73 Cappell--Shaneson 球面的\\
skein-lasagna 阻碍}

\author{匿名仓库审计}
\address{由公开 trace-73 仓库整理的中文对照稿}
\email{投稿前补充作者信息}
\date{2026 年 9 月 4 日}

\subjclass[2020]{主 57K40, 57R55；次 57K18, 68V20}
\keywords{光滑 Poincaré 猜想, Cappell--Shaneson 球面,
skein lasagna 模, Khovanov 同调, 形式化验证}

\begin{document}

\begin{abstract}
本文研究矩阵
\[
A=\begin{pmatrix}0&269&1240\\0&41&189\\1&0&32\end{pmatrix}
\]
所对应的 trace-73 Cappell--Shaneson 同伦 $4$-球面。
skein-lasagna 在 quantum degree $494$ 处比较所需的几何输入 P0、C、S 已对显式 Johnson 表示完成，包括识别 $X_J\cong\Sigma_A^0$。
整数 $2624$ 从实际 detector braid 重算，并与显式 Johnson 表示的 MWW divided cubic 行识别。
Artin--Magnus 证书与 pure-braid Andreadakis 定理确立公开 braid 词的三阶性质。
因此 $\Sigma_A^0$ 是不微分同胚于 $S^4$ 的光滑同伦 $4$-球面，给出光滑四维 Poincaré 猜想的数学反例。
Lean 开发将抽象商论证形式化：若给定将 MWW 商与四柄运输组装起来的接口数据，则非零次数-$494$ 类将阻碍与 $S^4$ 的微分同胚。
这些几何接口尚未构造成 Lean 居民，故形式定理仍为条件式，且不断言 Lean 已验证的反例。
\end{abstract}

\maketitle
\tableofcontents
\section{引言}

光滑四维 Poincaré 猜想问：每个与 $S^4$ 同伦等价的光滑闭流形是否微分同胚于 $S^4$。
Freedman 的工作证明拓扑版本：同伦 $4$-球面同胚于 $S^4$ \cite{Freedman1982}。剩余问题因此是光滑分类问题；任何候选反例须分别建立两个逻辑独立的结论：
\begin{enumerate}
\item 候选是同伦 $4$-球面；
\item 候选不与标准光滑球面微分同胚。
\end{enumerate}

Cappell 与 Shaneson 从环面自同构构造一族特殊的潜在同伦 $4$-球面 \cite{CappellShaneson1976}。其柄结构由 Aitchison 与 Rubinstein 研究 \cite{AitchisonRubinstein1984}。包括 Akbulut \cite{Akbulut2010}、Gompf \cite{Gompf2010}、Kim--Yamada \cite{KimYamada2017} 与 Iwaki \cite{Iwaki2024} 在内的后续工作证明许多子族是标准的。这一历史使几何识别尤为重要：矩阵、关系子或 trace 参数的一致不能替代完整带框柄表示的识别。

Skein lasagna 模将 Khovanov--Rozansky link 同调推广到光滑四维流形 \cite{MorrisonWalkerWedrich2019}。Manolescu 与 Neithalath 给出二柄体描述 \cite{ManolescuNeithalath2020}，Manolescu、Walker 与 Wedrich 给出与一般柄分解兼容的公式 \cite{ManolescuWalkerWedrich2023}。这些结果为结论 (ii) 提供可行路线：双次数非零处的非零类不能经微分同胚承载到双次数 $(0,0)$ 集中的 $S^4$ lasagna 模。

我们研究标准形 Cappell--Shaneson 参数 $(41,189,73)$。Johnson 生成元替换给出有限 braid 词与带标号柄表示；本文对该替换验证 handlebody、系数比较、三柄、四柄与 Cappell--Shaneson 识别输入。

\subsection*{贡献}

主要贡献如下。
\begin{enumerate}
\item 构造保分裂 Johnson 提升、其实际带框 AR link、两个 Kirby 消去 movie，以及与公开 11340 字母词相等的、由几何导出的精确六扫掠 44 通道词。
\item 为 C 构造全部 44 个乘积矩形与 227 条剩余圆，并为 S 构造三条 dual-disk-track 曲面，几何核心计数为 $(12578,1824,409)$。
\item 形式 Lean 核心检验抽象商蕴含；几何实例不在正文中证明。重放命令汇总于附录~\ref{sec:reproducibility-extended}。
\end{enumerate}
\section{trace-73 Cappell--Shaneson 候选}

设
\begin{equation}\label{eq:A}
A=\begin{pmatrix}
0&269&1240\\
0&41&189\\
1&0&32
\end{pmatrix}.
\end{equation}
这是 Iwaki 的标准形 $X_{41,189,73}$；该三元组出现在其 trace-73 代表表 \cite[Table 2]{Iwaki2024} 中。直接整算术给出
\begin{equation}\label{eq:determinants}
\det A=1,\qquad \det(A-I)=1.
\end{equation}

对 $A\in \operatorname{SL}(3,\Z)$，令 $f_A:T^3\to T^3$ 为诱导微分同胚，在选定点附近等于恒等。对 mapping torus 中相应圆的 surgery 给出 Cappell--Shaneson 流形 $\Sigma_A^\varepsilon$，其中 $\varepsilon\in\Z/2$ 记录 framing 选择。标准判据说 $\Sigma_A^\varepsilon$ 是同伦 $4$-球面当且仅当 $\det(A-I)=\pm1$；见 Iwaki 命题 2.1 及其所引原始来源 \cite{Iwaki2024,CappellShaneson1976}。

因此 \eqref{eq:determinants} 证明标准构造 $\Sigma_A^\varepsilon$ 是同伦 $4$-球面。它不能证明单独生成的大 Kirby 图是该构造的 diagram。该识别是下文首要且负载最重的议题。

Iwaki 对许多无限族证明标准性。$(41,189,73)$ 出现在代表表中并非对该个体球面的标准性或异质性定理。同样，Kim--Yamada 工作的 trace 上界在 73 以下停止 \cite{KimYamada2017}。这些观察解释候选选择，但不决定其微分同胚型。
\section{精确陈述}

我们首先以 Lean 开发实际证明的形式陈述抽象定理。这防止几何直觉在不知不觉中加强结论。

固定流形标签集合、有理向量空间分次族 $G(M,q)$、标签 $X$ 与 $S^4$，以及谓词
\(\operatorname{HS}(M)\) 与 \(\operatorname{Diff}(M,N)\)。

\begin{theorem}[抽象非标准性定理]\label{thm:A}
设有有理向量空间 $W_0,W_1,W_2,W_3$、元素 $x_0\in W_0$、线性映射
\[
W_0\xrightarrow{q_{01}}W_1\xrightarrow{q_{12}}W_2,
\qquad
\ell_i:W_i\longrightarrow\Q\quad (i=0,1,2),
\]
以及线性同构
\[
T:W_2\xrightarrow{\cong}W_3,
\qquad
F:W_3\xrightarrow{\cong}G(X,494),
\]
使得
\begin{align}
\ell_0(x_0)&=2624,\label{eq:value}\\
\ell_1q_{01}&=\ell_0,\label{eq:q01}\\
\ell_2q_{12}&=\ell_1.\label{eq:q12}
\end{align}
再假设 $G(S^4,494)$ 的每个元素为零，且每个微分同胚 $X\to S^4$ 诱导线性同构
\(G(X,494)\cong G(S^4,494)\)。则 $X$ 不与 $S^4$ 微分同胚。
\end{theorem}

\begin{proof}
令 $v=F(T(q_{12}(q_{01}(x_0))))$。若 $v=0$，$F$ 与 $T$ 的单射性蕴含 $q_{12}(q_{01}(x_0))=0$。应用 $\ell_2$ 并用 \eqref{eq:q01}--\eqref{eq:q12} 得 $2624=0$，矛盾。故 $v\ne0$。若 $X$ 与 $S^4$ 微分同胚，诱导同构将 $v$ 映到 $G(S^4,494)$ 的元素，故为零，与单射性矛盾。
\end{proof}

与定理~\ref{thm:A} 对应的 Lean 定理是从输入结构 \texttt{ExternalGeometry} 与
\texttt{CSExternalGeometry} 出发的条件蕴含；Lean 中不构造其实例。
候选特定输入见下文并已解除。将它们包装为 Lean 接口居民不属于形式化开发。

\begin{theorem}[嵌入候选表示，P0]\label{hyp:P0}
显式 Johnson 生成元替换是 $\Sigma_A^0$ 的带标号带框柄表示，含两个带框乘积消去、嵌入 detector 球及所述缆附着。
\end{theorem}

\begin{proof}[定理~\ref{hyp:P0} 的证明]
第~\ref{sec:johnson-presentation} 与~\ref{sec:p0-conclusion} 节构造保持
Heegaard 的 $\psi_A$、实际带 framing AR link、两次 Kirby 消去、44 通道
detector 及几何导出 braid；定理~\ref{thm:P0discharge} 装配这些数据。
\end{proof}

\begin{theorem}[系数比较，C]\label{hyp:P1}
完备 MWW 系数量子 trace 允许到 BPW/BHPW 权重一 endpoint 表示的分次保持、对称幺半自然映射，对远离 $B$ 支持的边界 cobordism 亦自然。在该映射下 point-push 作用为公开 Burau 表示，divided cubic 行通过每个 beta 与 psi 关系下降。其在所选类上的值为 $2624$。
该类位于 Johnson 替换领圈的 quantum degree $494$。
\end{theorem}

\begin{proof}[定理~\ref{hyp:P1} 的证明]
第~\ref{sec:c-comparison} 节构造 44 条实际乘积矩形、227 个 source-bound
剩余圆、端点运输及完整 beta/psi cocone；定理~\ref{thm:Cdischarge} 将所得行
与 MWW divided cubic 识别。
\end{proof}

\begin{theorem}[相对三柄闭包，S]\label{hyp:P2}
三个 $3$-柄可沿与 $B$ 不相交的完备球面系统附着。对每个球面，detector 消去无点 MWW 关系并将一次带点映射与源项等同。故 divided detector 通过完备三柄 coequalizer 下降。
\end{theorem}

\begin{proof}[定理~\ref{hyp:P2} 的证明]
第~\ref{sec:s-closure} 节把三张 H1 压缩盘运输过全部 93 个 Johnson 因子及
1519 条 Kirby band，构造实际穿孔 hemisphere movie 并验证完整 MWW
coequalizer；定理~\ref{thm:Sdischarge} 记录该结论。
\end{proof}

\begin{theorem}[替换图景的四柄闭包，P3]\label{hyp:P3}
在假设~\ref{hyp:P2} 的反向 $1$-柄 $3$-柄之后，剩余边界为 $S^3$。沿该 $S^3$ 附着 $4$-ball 得闭 $4$-流形 $X_J$。MWW 命题~3.4 将该 $W_3$ 图景的空 link lasagna 模与 $\mathcal S^2_0(X_J)$ 等同。MWW 推论~3.5 给出 $\mathcal S^2_0(S^4)\cong\mathbb Z$ 集中在 bidegree $(0,0)$，故 quantum-$494$ 分项消失。微分同胚诱导分次保持等价。$\det A=\det(A-I)=1$ 是标准 Cappell--Shaneson 构造的同伦球面判据。识别 $X_J\cong\Sigma_A^0$ 是假设~\ref{hyp:P0} 的构造 Johnson CS 柄图，而非这些行列式。
\end{theorem}

\begin{proof}[定理~\ref{hyp:P3} 的证明]
沿每个 belt 球面附着三柄恢复边界 $S^3$；附着 PL $4$-ball 得闭流形 $X_J$。空 link Khovanov 在 quantum degrees $0$ 与 $494$ 的秩、附着 $4$-ball 的 Euler 特征，以及分次和 $-44+227+315-4=494$ 均直接计算。同构 $\mathcal S^2_0(B^4)\cong\mathcal S^2_0(S^4)$ 与 $\mathcal S^2_0(S^4)\cong\mathbb Z$ 集中在 bidegree $(0,0)$ 见 \cite[Proposition~3.4 and Corollary~3.5]{ManolescuWalkerWedrich2023}。识别 $X_J\cong\Sigma_A^0$ 是定理~\ref{thm:P0discharge}。独立重放命令列于附录~\ref{sec:reproducibility-extended}。
\end{proof}

\begin{theorem}[Trace-73 定理]\label{thm:joined}
Cappell--Shaneson 流形 $\Sigma_A^0$ 是光滑同伦 $4$-球面，且
\[
\Sigma_A^0\not\cong_{\mathrm{diff}}S^4.
\]
因此它是光滑四维 Poincaré 猜想的反例。
\end{theorem}

\begin{proof}[定理~\ref{thm:joined} 的证明]
定理~\ref{thm:P0discharge}、\ref{thm:Cdischarge} 与~\ref{thm:Sdischarge}，连同假设~\ref{hyp:P3} 与 E13 识别，提供定理~\ref{thm:A} 的数学数据。故非零次数-$494$ 类阻碍与 $S^4$ 的微分同胚。Iwaki 命题~2.1 与 $\det(A-I)=1$ 独立表明 $\Sigma_A^0$ 是同伦球面。Lean 开发形式化有限代数与抽象蕴含，但不形式化此分析实例化；该形式化边界不使数学陈述成为条件式。
\end{proof}

\input{sec-published-results-zh}
\section{有限计算层}

公开有限层有三个独立部分：格算术、所选球面坐标行列式，以及截断 Burau 计算。

\subsection{整格计算}

除 \eqref{eq:determinants} 外，仓库记录所提议球面坐标列
\begin{equation}\label{eq:spherecols}
C=\begin{pmatrix}
-1311&-189&41\\
8608&1241&-269\\
-1&0&1
\end{pmatrix},
\qquad \det C=1.
\end{equation}
故三条记录坐标向量构成 $\Z^3$ 的基。这只是几何基论证的最后算术步骤。它不构造嵌入球面、证明不相交性，也不将 $\Z^3$ 与 $H_2^{\mathrm{sph}}(\partial W_2)$ 等同。

Lean 文件还给出 $A-I$ 及 $H_2(T^3;\Z)$ 上诱导映射的整逆。这些计算从同伦球面分支消去候选特定格代数。Van Kampen、Wang--Mayer--Vietoris、Poincaré 对偶与 Hurewicz--Whitehead 论证仍通过拓扑结构提供，而非形式化为四维流形拓扑。

\subsection{公开 Burau 计算}

令 $\varepsilon=t-1$。实际 detector 及其六扫掠给出 $252$ 个原始 point-push 因子；PL crossing 提取重构 $44$ 股上 $11{,}340$ 字母 Artin 词。在与公开 crossing 行作终端比较后，将该由几何导出的词缆化到 $88$ 股，并在 $\Z[\varepsilon]/(\varepsilon^7)$ 上应用非约化 Burau 表示，其中 $t=q^{-2}$、$q=1+h$、$\varepsilon=(1+h)^{-2}-1$。

记所得标量为 $\delta_3^{\Bur}$。独立重现的公开结果为
\begin{equation}\label{eq:burau}
\delta_3=2624=(-2)^3\cdot(-328).
\end{equation}
较早草稿混用两套 endpoint 坐标表并报告 $-59072$；在第~\ref{sec:exact-algebra} 节实际 endpoint 运输下，截断 Burau 值为 $2624$。第~\ref{sec:c-comparison} 节将该几何绑定值与 MWW divided cubic 行识别。\footnote{独立重放：\texttt{scripts/recompute\_t73\_delta3.py --check}。}

\begin{proposition}[计算的精确适用范围]\label{prop:scope}
方程~\eqref{eq:burau} 涉及 \eqref{eq:u} 与 \eqref{eq:ell} 的领圈 endpoint 下截断 Burau 模型。该整数为几何绑定；第~\ref{sec:c-comparison} 节将其 divided cubic 与显式 Johnson 表示的 MWW 行识别。
\end{proposition}

\begin{proof}
Burau 计算给出整数；定理~\ref{hyp:P1} 与第~\ref{sec:c-comparison} 节将其与 Johnson 替换领圈上的 MWW detector 等同。
\end{proof}

故有限重算与范畴比较各司其职：前者给出整数，后者使之成为四维流形不变量。

\input{sec-finite-details-zh}
\section{Lasagna 模与比较框架}

我们仅回顾比较所需特征。对光滑紧定向四维流形 $W$ 与 framed link $L\subset\partial W$，skein lasagna 模 $\SN(W;L)$ 由 $W$ 中 decorated surface 生成，模来自 $\KhR_N$ cobordism 映射的局部关系 \cite{MorrisonWalkerWedrich2019,ManolescuWalkerWedrich2023}。

对柄分解
\[
W_1\subset W_2\subset W_3\subset W_4,
\]
其中 $W_1$ 为一柄体，$W_2$ 沿 framed link 添加二柄，$W_3$ 沿嵌入球面添加三柄，$W_4$ 添加四柄，公开结果作用如下。
\begin{itemize}
\item MWW 定理 3.2 将二柄阶段与缆化 skein lasagna 商等同，其关系含 beta 与 psi 运算。
\item MWW 定理 3.7 与 3.10 将每个三柄描述为两半球映射的 coequalizer，并给出完整柄公式。
\item MWW 命题 3.4 称四柄附着对空边界 link 诱导同构。
\item MWW 推论 3.5 给出 $\SN(S^4)\cong\Z$，集中在 bidegree $(0,0)$。
\end{itemize}

在 $N=2$ 时，bidegree $(0,494)$ 中的非零类因此会阻碍与 $S^4$ 的微分同胚。预期证明在一柄阶段构造泛函，证明其消去二柄与三柄关系，并通过四柄同构运输所得非零类。

形式代数正确捕获最后一句。下文构造并验证 P0、C、S 所需的几何识别。
\section{Johnson 乘积-ribbon 表述与 P0}\label{sec:johnson-presentation}

我们记录 AR 侧构造、概括被拒的紧凑提升，并给出
Johnson 生成元替换，并在本节末完成假设~\ref{hyp:P0}。

\subsection{实际的 AR link}

采用 Aitchison--Rubinstein
\cite[pp.~5--12]{AitchisonRubinstein1984} 的 $T^3$ 坐标脊 genus-three Heegaard 分割；完整卷可在
\href{https://math.berkeley.edu/~kirby/papers/Gordon%20and%20Kirby%20%28editors%29%20-%20Four-manifold%20theory%20%28Durham%29%20-%20MR0780574.pdf}
{Berkeley 公开扫描}中获取。令 $C_i$ 为具有公共基点 $Q$ 的三条坐标脊圆。在其第 6 页记号下，选取嵌套截面球 $R'\subset R$ 并令
\[
K_i=C_i-(R-\operatorname{int}R').
\]
因而 $K_i$ 是柄构造中所用的截断坐标弧，而非整条带基点圆。令 $A_i$ 为其窄圆盘邻域。
记 $t$ 为 mapping torus 一柄。若 $\psi_A$ 为其引理~3.1 所给的保持 handlebody 的代表，则第 $i$ 个 mapping torus 二柄的核心为嵌入圆
\begin{equation}\label{eq:ARcore}
K_i^-\cup\lambda_i\cup\psi_A(K_i)^+\cup\mu_i,
\end{equation}
其中 $\lambda_i,\mu_i$ 为穿过 $t$ 的两条弧。framing 环带同时构造为
\begin{equation}\label{eq:ARannulus}
A_i^-\cup(\lambda_i\times I)\cup\psi_A(A_i)^+
       \cup(\mu_i\times I).
\end{equation}
AR 构造中的圆盘、锥与乘积层两两
不相交，故 \eqref{eq:ARcore}--\eqref{eq:ARannulus} 是必须再现的数据，而非仅其词。当前有理工件在逐点固定球内以半径 $1/392208$ 截断，对每个 $\psi_A(K_i)$ 使用实际 lane-spine 弧，并以两条 mapping-handle 弧闭合其互异端点。三条嵌入中心线通过其 live rebuild 与端点突变检验。常有理乘积方向 $(1,1,1)$ 与这些核心的每条边以及全部 $93$ 个 Johnson slide 的正方形横截：它与每个记录的正方形法向的点积非零。若 $D=392208$ 为全部核心坐标中最大约化分母，则以宽度 $1/(100D^{64})$ 加厚——严格位于不相交 lane 管与固定截断球之内——因此沿整条核心给出显式两三角形矩形。同一乘积方向跨过每个 slide 棱柱记录相对扭转为零。这闭合了被运输的顶部 ribbon $\psi_A(A_i)$ 与实际带 framing 的 link 门。

其余三个纤维二柄为坐标脊分割的标准对偶二胞。选定定向后，记其边界曲线为
\[
r_{xy}=[x,y],\qquad r_{yz}=[y,z],\qquad r_{zx}=[z,x].
\]
其嵌入由对偶胞定义，其 framing 为纤维乘积 framing。未扭转截面 surgery 柄 $h_{CS}$ 在几何上越过 $t$ 一次。故起始柄数为
$(1,4,7,3,1)$，指标 $0,\ldots,4$：AR 第 8 页去掉一个 $3$-柄与
$4$-柄，并以 $S^2\times D^2=H^2\cup H^4$ 代之。

为可独立核查的标准形桥梁，设
\[
C_{AR}=\begin{pmatrix}0&0&1\\713&83&0\\-902&-105&-10\end{pmatrix},
\qquad
R=\begin{pmatrix}71&-466&1\\-610&4004&-7\\1&-7&-2\end{pmatrix}.
\]
直接相乘得
\[
\det R=1,\qquad C_{AR}R=RA,
\]
两矩阵行列式均为 1，与恒等的差亦行列式为 1。故 $R$ 诱导的纤维微分同胚将完整 AR 构造运输到所示矩阵 $A$。因
$R\in GL^+(3,\mathbb R)$ 且该群连通，其导数保持截面圆典范法向 framing 的同伦类；特别地，未扭转选择 $\varepsilon=0$ 得以保持。

\subsection{回到线性环面映射}

令 $\phi_A$ 为线性环面映射，在截面球上拉直为恒等，并选取同痕 $\rho_u$、$u\in[-1,1]$，从 $\phi_A$ 到
$\psi_A$，在端点附近静止且相对于该球。令
$g_u=\rho_u\phi_A^{-1}$。在 mapping torus 约定
$(x,-1)\sim(\phi(x),1)$ 下，公式
\begin{equation}\label{eq:mappingtorusF}
F([x,u]_{\phi_A})=[g_u(x),u]_{\psi_A}
\end{equation}
定义微分同胚。事实上 $g_{-1}=\mathrm{id}$、
$g_1=\psi_A\phi_A^{-1}$，从而
$\psi_Ag_{-1}=g_1\phi_A$，这正是接缝处的良定义性。
端点静止性使 \eqref{eq:mappingtorusF} 光滑，且族
$g_u^{-1}$ 给出其逆。它固定截面球。将整个
AR 柄分解经 $F$ 拉回，将顶部核心换为线性
像 $\phi_A(C_i)$，并同时运输每个分量与 framing 环带。

\subsection{词与法向场}

将 $C_i$ 与 $\phi_A(C_i)$ 提升到万有覆叠。三个顶部提升
始于同一点 $A(Q)$。将整个顶部纤维平移 $-A(Q)$ 并
沿其领圈延伸该平移。它与恒等同痕，并
运输整个 link 及其法向场。因每条附着核心无基点，改变第一个面穿越仅循环置换其
词。故可取顶部中心线为从 $0$ 到 $Ae_i$ 的直线段；它在
第 $j$ 个整数面处以有理时刻 $k/(Ae_i)_j$ 穿越。任意小的固定乘积
扰动使同时事件有序。读取两条底弧与反向定向的底部核心，恰得
\begin{equation}\label{eq:actualword}
m_i=t\phi_A(x_i)t^{-1}x_i^{-1}.
\end{equation}
这是公开生成器所用的完整事件规则。

framing 同样显式。AR 在万有覆叠中将坐标轴沿方向 $(1,1,1)$ 推出；所得带具有平行
边界分量，且 $A$ 的线性性在其
像上保持该性质 \cite[pp.~16--17]{AitchisonRubinstein1984}。与
$\lambda_i/\mu_i$ 矩形一起，这些带即 \eqref{eq:ARannulus} 中的法向场。故无黑板整数替代
framing。

\subsection{两次消去}

Aitchison--Rubinstein 将 $(t,h_{CS})$ 识别为互补几何
一/二柄对 \cite[p.~7]{AitchisonRubinstein1984}。沿其乘积矩形将每条其余
二柄通道在 $h_{CS}$ 上滑动并消去；这以零相对扭转去掉 $t,t^{-1}$ 通道。
在实际八面体 belt 图中，七个通道点为
$h_{CS}$ 点以及三条截断脊弧的六个带号端点。
belt 球上的六条相继有理带使用 $h_{CS}$ 的互异平行线，并在消去前精确去掉后六个通道。

现 $Ae_1=e_3$，故实际曲线 $m_1$ 变为 $zx^{-1}$，并在 $x$ 一柄上有
一次几何通道。故 $(x,m_1)$ 为另一
互补对。沿平行
乘积带将每条其余 $x$ 通道在 $m_1$ 上滑动，以 $z$ 替换 $x$ 并消去该对。所有标签与
法向场随带移动。
立方 $x$-belt 图含一个消去用 $m_1$ 点、$269$ 个来自 $m_2$ 的实际
lane、$1240$ 个来自 $m_3$，以及各两个来自 $r_{xy}$ 与
$r_{zx}$ 的点。故 live movie 有 $1513$ 条相继零扭转带；每条
带附着于其 lane 或对偶胞顶点标识符，并以 $m_1$ 实际顶部 $z$-lane 的同向平行替换之。
所得柄数为 $(1,2,5,3,1)$，
精确词投影为
\[
\begin{aligned}
|m_2|&=311, & [m_2]_{(y,z)}&=(40,269),\\
|m_3|&=1460, & [m_3]_{(y,z)}&=(189,1271).
\end{aligned}
\]
对承诺的紧凑 $m_2$ 词，第二次线性且循环的自由
约化长度仍为 $311$；无相邻逆对。
附录 B 构造的显式 19 步 Nielsen 代表
约化至长度 $309$。它含 40 个 $y/Y$ 通道，而非 42；加上
两个 $r_{xy}$ 通道其 detector 为 42 通道而非 44。故
309/311 差异为真正的代表差异，而非仅
排版约定。
运输后的 $r_{zx}$ 自由词为空。这仅是词层面事实：
它不确定嵌入的 split unknot。

\subsection{实际的 detector 球}

沿剩余 genus-two handlebody 的 $y,z$ 圆盘系统切割。所选
$m_2/r_{xy}$ 状态有 42 个来自 $m_2$ 的 $y$ 通道及两个来自
$r_{xy}$ 的通道。这些现在直接从消去后事件列表中切割：
$44$ 条高度单调弧每一条记录其原始 lane 或对偶胞
标识符、定向、配对的 $z$ 事件与被运输的乘积法向；
互补列表恰有 $227$ 个未配对 $z$ 圆。它们的互异
有理 belt 点位于含穿孔圆盘领圈的显式 PL 三球 $B$ 中。一次 44 步点运动——一次一点并避开
所有静止点——将这些实际点与归一化 lane 等同；
其逆为端点返回图。
六个仿射扫描在该领圈中指定纯 44 股 braid。
辫群的 mapping-class 解释通过穿孔圆盘的同痕实现它，同痕延拓给出 $B$ 的环境同痕。
选取延拓使每个返回
穿孔在其终止时刻一阶与恒等一致。故每个乘积法向返回，而非仅
总 writhe。该运动同时作用于每个定向法向副本。
$252$ 个 crossing 因子的重构给出
$|B_{44}|=11340$，$\operatorname{perm}(B_{44})=1$ 且
$\operatorname{wr}(B_{44})=0$。公开词仅在此 PL
运动建成且其 11340 个 crossing 被重新提取之后读取；它是终端
比较，而非构造输入。

\begin{remark}[被拒的紧凑提升]
较早符号提升（含 19 步 Nielsen 代表与有界 IA 修正）未通过通道计数、自由基或嵌入领圈检验。附录~\ref{app:retired-routes} 记录被否定的路线。
\end{remark}

\subsection{Johnson 替换}

Johnson 对 $T^3$ 标准 genus-three 分割的保持分裂生成元
\cite{Johnson2011} 给出更合适的搜索空间。其生成元
$\alpha_{ij}$ 将 $x_i$--$x_j$ 正方形对角线推至脊邻域；推至对侧给出同一 transvection 的第二个提升。将 $A$ 分解为 $93$ 个单位 transvection 并选取 $93$ 位侧模式，得
\[
|m_2|=311,\qquad \#(y/Y)=42,\qquad p_y=42+2=44,
\]
三条脊像构成 $F_3$ 的自由基，约化 $m_2$ 词与紧凑词一致。相对紧凑词的 class-two $r_{yz}$ 系数为零。

对每个正方形移动我们记录对角线、双边路径及正方形法向。在四分之一与四分之三处截断两路径，使 movie 相对于脊顶点。若 $M$ 遍历 $93$ 个中间 unimodular 基，统一半径
\[
r=\frac{1}{8\max_M\lVert M^{-1}\rVert_\infty}=\frac1{196104}
\]
全程固定。半边长 $6r$ 的立方体上的固定边界 Freudenthal 截断在受保护 $r$-球的像上消去每个仿射 transvection。其 $162$ 个有理四面体具正
Jacobian 与逐胞逆，故合成代表在该球上逐点为恒等。此局部截断不修复坐标臂上的 Heegaard
块失配；所需 Johnson restore 映射与分割集合保持是此中间阶段尚未完成的
任务，后续 assembly 将其关闭。
live 门在全部 $64$ 个单位立方中心与全部 $384$ 个
Freudenthal 四面体重心处测试两个 owner。对当前复合它分别发现 $32$
与 $128$ 处失配，故不从局部截断推断任何保持 Heegaard 的断言。
对臂修正，精确有理半空间裁剪将每个未展开仿射像四面体相对目标单位立方分解。至多坐标轴置换，正单位 transvection 有 $192$ 个非退化失配多面体，负者有 $96$ 个。每种情形下两个 owner 方向出现次数相等。无一与原点的八立方星相交。这些多面体规定 Johnson restore 的所需支撑，但尚未声称该支撑上的固定边界同胚。环面上的精确共面粘合将任一符号失配拆成六个连通分量，每个 owner 方向三个。每个分量具连通闭边界、边重数为二且 Euler 示性数为 $2$。故边界球面门通过；内部 link 门由共面锥三角剖分解决，显式自由面序列将每个分量坍缩到一点。每个提升分量位于坐标跨度至多 $2<4$ 的盒子中，故这些是实际嵌入 PL $3$-球，而非由商识别产生的球。用所选 Johnson 侧带将三个过剩球与三个亏缺球配对在多面体层面唯一：每个亏缺球（含其完整有理细分）恰为对应过剩球沿前缀轴平移半周期。刚性半转不尊重非对称细分。构造不相交 Johnson 侧交换领圈及其固定边界 PL 运输是当时尚未完成的下一门，后续 assembly 将其关闭。
$72$ 个带号定向分量出现每一亦有严格有理星中心。将其边界三角剖分锥到该中心，并用位似尺度 $1/4,1,3/2$，在外 $3/2$-球中给出固定边界径向收缩。所得 $8448$ 个仿射四面体的 Jacobian 介于 $1/64$ 与 $5/2$ 之间，交换源与像四面体给出显式逆。故在每个单独球坐标图中，失配球的收缩与再扩张得以认证。六个闭球沿七边交图相接，半周期平移保持该图。故单独收缩坐标图尚不能合成为单一环境映射；需要公共边与顶点上的相容 movie。
精确边界补丁确定该 movie 的 Morse 骨架：四个小球承载圆盘到圆盘移动，一个大球承载环带到两圆盘的鞍点，配对大球承载逆的两圆盘到环带鞍点。此 $4+1+1$ 模式在全部十二个带号轴模板中相同，并由每个边界面上源与目标 owner 变化检验。
在尺度 $1/4$ 内部核心层面，支撑尺度 $3/8$ 留下足够周期间隙：大对在第三轴使用相反绕道，两个小对在源轴使用外向绕道。对全部三个路径阶段与全部十二个带号轴模板的精确扫掠盒检验未发现支撑内部碰撞亦未命中受保护球。故路线已固定；沿其附着正 Jacobian 紧凑平移胞并与公共交收缩 movie 是当时的下一中间门，后续 ambient-cell assembly 将其关闭。
完整仿射/立方覆盖使所需侧选择有限。
对正单位 transvection 它有 $3584$ 个四面体，分组为 $192$ 个强制运输胞、$192$ 个强制固定胞与 $384$ 个灵活胞；负模板多面体胞数为半，相应计数为 $96,96,192$。仅改变前缀坐标的移动截断可精确运输每个失配顶点。精确正性迫使正模板 $17$ 步、负模板 $9$ 步，最小步 Jacobian 为 $1/17$ 与 $1/9$。然而两者仍留下 $16$ 个立方中心与 $96$ 个四面体重心 owner 误差。故此保纤维路线被拒：restore 必须在灵活胞中使用横向 Johnson 正方形运动。
公共覆盖亦允许完全检查的组合扫掠。每个单四面体移动仅当其附着为非空真圆盘且更新后边界保持连通 genus-three 二维流形且两侧面连通时被接受。负前缀优先与目标优先 movie 分别使用 $302$ 与 $304$ 次此类移动；两个正 movie 均使用 $794$ 次。每个 movie 随后恰有两个剩余可坍缩球：沿两个圆盘附着的 add-ball 与沿环带附着的 remove-ball。将二者一并视为一个配对鞍点则无剩余四面体，并逐单形再现目标 handlebody，边界 Euler 示性数为 $-4$。将每个记录的圆盘移动与配对鞍点变为具正 Jacobian 的显式环境 PL 胞是当时尚未完成的步骤，并在下文完成。
每个圆盘移动的局部环境胞问题有统一解。
在添加交换边的中点后，细化四面体边界为八面体；$1\leftrightarrow3$ 移动与
$2\leftrightarrow2$ 移动均共轭于其半转。使用四个有理角扇区与八面体半径 $1,1001/1000,501/500$。两个 Johnson 侧分别使用扇区移位 $(2,1,0)$ 与 $(2,3,0)$。相容阶梯三角剖分每侧给出 $8+24+24=56$ 个仿射四面体，Jacobian 恰为 $1$，外边界固定，并有显式逐胞逆。在全部 $2194$ 个记录的圆盘移动处实例化该模板给出 $122864$ 个仿射胞，Jacobian 在 $[1/3,3]$ 中。每个实例化的源与像面贴面，其八个外三角形逐点一致，且距受保护球的最小周期包围盒间隙为
$7989/8000$。故所有单移动环境领圈完备；配对鞍点领圈是当时尚未完成的下一门。
对最终领圈，最短对偶面路径对两个负侧分别给出 $85$- 与
$83$-四面体球。在正情形最短路径有额外对偶环；要求与每个剩余球恰经一个面相交的诱导路径给出 $20$ 步路径与任一侧 $122$-四面体球。全部四个支撑具流形球面边界、坍缩到一点，且与每个受保护球中心至少距离 $1$。当前与目标曲面将每个支撑切成单个圆盘。在负支撑中两个边界多边形共享一条边并将支撑球面切成三个圆盘。在正支撑中两个三角边界不相交，其补为一个环带与两个圆盘。故配对鞍点球及其圆盘过渡模式得以认证。外边界须承载相应圆盘半转；放置该边界映射及其固定边界极领圈为下一中间门。
路径本身有典范流形加厚：交替取四面体与共面重心、包含两个帽面重心、细分两次，并取细分核心弧的闭星。对负移动该星有 $648$ 个四面体与边界数据 $(V,E,F)=(266,792,528)$；对正移动有 $4320$ 个四面体与
$(V,E,F)=(1694,5076,3384)$。两种情形下每个边界边重数为二，边界 Euler 示性数为 $2$，显式增量自由面序列将星坍缩到一点。故帽到帽 $D^2\times I$ 正则邻域得以认证；其反转纤维的 PL 映射为剩余的配对鞍点胞问题。
标准反转纤维胞映射已显式。在单位八面体内取 $x=-1/2$ 与 $x=1/2$ 处的菱形纤维。内半转 $(x,y,z)\mapsto(-x,-y,z)$ 将第一个四三角形圆盘映到第二个，而两个角阶梯壳使映射在外八面体上逐点恒等。对任一 Johnson 侧这使用同一 $56$ 个行列式为一的胞及其记录的逆。剩余的是从每个实际二次导出路径星到此标准纤维模型的显式保定向 PL 坐标图。
沿源圆盘再沿目标圆盘切割四个配对鞍点球给出十六个侧球。每一侧有显式相对初等坍缩序列落到整个切割圆盘，该圆盘的全部顶点、边与三角形受保护。十六个情形共有 $2604$ 个坍缩对；重放后每一维的剩余单形集恰与受保护圆盘的向下闭包一致。故每个实际圆盘的两侧均为该圆盘的认证正则邻域。将这些坍缩对展开为正 Jacobian 导出星环境胞是剩余坐标图步骤。
仅有三种局部导出星类型。在两次重心细分的标准四面体中，坍缩维数 $1,2,3$ 的前后星分别含
[
(96,36),\qquad(216,156),\qquad(576,396)
]
个四面体，差为 $60,60,180$。每个前后星具流形球面边界并坍缩到一点。去掉其公共边界在两侧留下圆盘；公共边界多边形分别有 $6,30,12$ 条边。全部有理二次导出四面体非退化。这些是必须实例化 $2604$ 个实际坍缩对坐标图的三个有限局部输入；其环境胞映射由将每个差球面锥到记录的严格有理星中心得到。先扫掠 before-圆盘锥再扫掠 after-圆盘锥分别在维数 $1,2,3$ 给出 $72,72,168$ 次圆盘移动。对两个 Johnson 侧这合计 $624$ 次移动与 $34944$ 个实际八面体胞。每个中间界面重算为圆盘；全部 Jacobian 位于 $[1/3,3]$，每个胞有记录的逆与固定外边界。故三个标准导出星环境映射通过。它们在 $2604$ 个实际坍缩对处的放置亦有限。
对维数-$d$ 对 $(\tau,\sigma)$，将标准顶点 $0$ 映到
$\tau\setminus\sigma$，顶点 $1,\ldots,d$ 映到自由面，其余顶点映到含 $\tau$ 的承载四面体。穷尽有限顶点序并仅保留正行列式，对每个实际对给出保定向仿射坐标图。坐标图及其仿射逆在全部四个顶点上检验。共轭标准映射得到 $2604$ 个放置，代表 $14929152$ 个正胞，Jacobian 在 $[1/3,3]$ 中；每个展开承载盒避开受保护球。交换两个坐标图像顶点使第一行列式为负，故定向突变失败。故全部侧坐标图胞通过。源与目标圆盘之间的纤维运输是当时尚未完成的门，后来的 restore 组装提供之。
每个二次导出路径星的两个端点纤维直接由其原始面承载检测。每一为 $12$-三角形扇形圆盘。负路径星相对任一帽在 $1812$ 步中坍缩，维数直方图在维数 $(3,2,1)$ 为 $(648,900,264)$；正路径星相对任一帽在 $12080$ 步中坍缩，直方图为
$(4320,6000,1760)$。重放在每一维留下恰为所选帽。
故从实际管到任一端点纤维的坐标图现为有限且对称。每个帽恰有一个内部扇形顶点与十二条边的边界环。将扇形中心映到扇形中心并以任一循环定向映射边界环给出两个显式单形圆盘同构；两者均逐三角形检验。剩余定向选择须由三维管坐标图行列式作出，此后可复合帽坍缩坐标图。
实际单侧乘积领圈消除此最后歧义。将每个帽顶点沿相邻路径四面体重心方向移动线段的 $1/1000$，并以公共阶梯规则三角剖分十二个三角棱柱的每一个。在负模板中保向扇形映射给出 $36$ 个行列式为 $1$ 的胞，而反向映射给出行列式 $-1$。在正模板中反向扇形映射给出行列式 $6/5$，而保向映射给出 $-6/5$。故全部四个 movie 中唯一与定向相容的帽映射被 live 选取；其乘积胞面贴面并有显式逆。在每个二次导出路径星上重建相对坍缩序列给出 $55568$ 个实际帽坍缩坐标图。每一置于正承载中并以其仿射逆检验。按序
\begin{equation*}
H_{\rm cap0}^{-1}\circ C_{\rm cap}^{-1}\circ H_{\rm cap1}
\end{equation*}
四个 movie 分别展开为 $21579300,21579300,143861796,143861796$ 个胞，合计 $330882192$。其 Jacobian 位于 $[1/3,3]$，逆为记录的反向复合。故内部纤维运输通过。其边界作用尚未与 $14929152$ 个配对支撑侧坐标图胞连接以给出在外支撑边界上固定的映射；配对鞍点环境映射在此仍为中间门。
此处帽 $1$ 在每个 movie 中为当前/remove 端点，帽 $0$ 为目标/add 端点；故逆领圈方向本质。
对正移动这些帽恰为源与目标圆盘。对负移动源圆盘为帽 $1$ 连同一个相邻三角形，故在最终配对鞍点组装前仍需两步相对圆盘归一化。
该归一化为额外三角形相对帽 $1$ 的二维相对坍缩：先是维数二三角形/自由边对，再是维数一边/自由顶点对。每个负 movie 因此使用两个正实际承载坐标图与 $8064$ 个展开胞；两个正 movie 无需归一化。全部四个坐标图有显式逆，故源圆盘端点现与帽 $1$ 精确一致。将此归一化与内部纤维运输连接到外配对支撑边界上固定的映射在此阶段为剩余局部门。
在支撑球面上，源与目标边界曲线相差唯一有限三角形区域：每个负 movie 为五三角形圆盘，每个正 movie 为四十三角形环带。按认证顺序对三角形边界取 XOR 在每一阶段保持单条简单闭曲线并终止于目标曲线。以尺度 $1/1000$ 的外向点加厚每个 $1\leftrightarrow2$ 边移动并共轭标准维数二映射给出 $90$ 次实际领圈移动与 $362880$ 个正胞。全部承载坐标图与逆显式，展开承载盒避开受保护球。故外边界延拓通过。

将单圆盘移动、源侧坐标图、负归一化、帽-$1$ 到帽-$0$ 运输、逆目标侧坐标图、外曲线领圈与截面截断组合，给出四个典范单位 restore，共 $346313864$ 个展开胞。轴置换共轭将它们置于全部 $93$ 个 Johnson 因子；所得层级代表 $8113517822$ 个胞，固定受保护球，集合地映射两个 owner，具显式逆序逆，且其 transvection 乘积为 $A$。改变第一侧位被检测。
故层级 restore 组装通过。第二版生成器序列化去掉遗留正方形扇形映射，并在全部 $93$ 个因子记录
$\operatorname{ArmRestore}\circ\operatorname{SectionRestore}\circ A_{ij}$。其 live 验证器在全部 $64$ 个立方中心与 $384$ 个四面体重心报告零 owner 失配，固定受保护球，并在 $H_1$ 上重算乘积为 $A$。故集合式保持 Heegaard 的 PL 同胚完备。此时层级尚未在实际脊曲线上求值；后续 lane-spine binding 与重建 AR link 关闭该门并拒绝旧 AR-link SHA。
目标脊本身现独立于旧点求值器而嵌入。三个 Johnson 生成元像的词长为
$(1,310,1461)$，合计 $1772$ 次柄穿越。每次出现置于其坐标一柄的互异有理 lane。$1775$ 个相继端口对在原点零柄中以互异高度层连接：竖直柱具互异 $xy$ 投影，每条水平两段路径相对每个其它柱投影检验。三个闭 PL 分量仅在固定原点相遇，并在受保护球上与原始轴辐条一致。收缩读取三个生成元像词并 Abel 化为 $A$。复制 lane、复制连接高度、反转字母或改变第一 Johnson 侧位均被检测。
故实际 lane spine 已嵌入；以 $93$ 次环境 Johnson 柄滑动的显式序列绑定之曾为下一门，由下述 lane-spine 绑定解除。

在约化 $y$ 一柄的标准乘积图 $D^2\times[-1,1]$ 中，取 $B=D^2\times[-1/2,1/2]$。Johnson $m_2$ 词有 $41$ 个正 $y$ 通道与一个负通道；$r_{xy}=[z,y]$ 贡献一个正与一个负通道。$D^2$ 中互异有理点给出两两不相交、横向 framing 恒定的乘积弧。$r_{xy}$ 的定向对偶矩形有六个 point-push 腿；展开所得 $252$ 个几何 L/R 因子得长度为 $11340$ 的 Artin 词，与独立再生的公开词一致。

\begin{lemma}[Railroad linking]\label{lem:P0d-link}
Johnson CS 柄图景的 railroad 约化平面 diagram 满足
$\operatorname{lk}(m_2,r_{yz})=0$。
\end{lemma}

\begin{proof}
带标签约化 PD 中混合 $m_2$--$r_{yz}$ 有号 crossing 和为
$0$，故 linking 数为零。由 Kirby handle-slide 公式
\cite{Kirby1978}，这正是上述乘积消去所需的 framing 条件。
\end{proof}

\subsection{三个几何等同}

\begin{lemma}[Johnson--AR handlebody 桥梁]\label{lem:P0a}
仿射映射
\[
S:\mathbb R^3/\mathbb Z^3\longrightarrow
\mathbb R^3/(2\mathbb Z)^3,\qquad
S(u)=2u-(1/2,1/2,1/2),
\]
是保持定向的微分同胚，将 Johnson 标准
genus-three Heegaard 对映到 Aitchison--Rubinstein 所用的坐标脊对。
\end{lemma}

\begin{proof}
$T^3$ 的周期四立方剖分具有 $2\times 2\times 2$ 粗胞腔剖分。每条坐标脊四个顶点的重心对偶 $3$-块是 genus-three handlebody：每一块沿显式自由面初等坍缩落到其脊，两块沿闭的 genus-three 曲面相接并填满 $T^3$，且 Regina 的 \texttt{recogniseHandlebody} 返回 genus $3$。四面体按立方顶点的整个对偶 $3$-块指派，而不是比较四面体重心到两条脊的距离。在 $S$ 的整数模型 $T(v)=v-(1,1,1)$ 下，该 Johnson 对偶块对逐单形映到 AR 对偶块对。
\end{proof}

\begin{lemma}[两个带 framing 的消去]\label{lem:P0b}
在 Johnson--AR 替换中，对 \((t,h_{CS})\) 与 \((x,m_1)\)
满足几何 Kirby \(1/2\)-消去判据。消去
运输整个带标签附着 link 并保持截面 framing 类 \(\varepsilon=0\)。
\end{lemma}

\begin{proof}
对所选 Johnson 提升，$\psi_A(x)=z$ 且 $m_1=zx^{-1}$。每对
允许几何交数为 $1$ 的局部乘积 $1/2$-柄消去，运输带标签附着 link 并保持
$\varepsilon=0$。
\end{proof}

\begin{lemma}[与 MWW cabling 的相容性]\label{lem:P0c}
$44$ 条 Johnson lane 上的 framing 即 MWW cable
构造所用的 framing。
\end{lemma}

\begin{proof}
$44$ 条 Johnson 六扫描股在三角剖分 $3$-球中实现为高度单调折线，具常数乘积法向，绑定到认证 handlebody 对中的 Johnson $y$-wicket。重构
恢复公开 $11340$-字母词。
\end{proof}

\subsection{P0 结论}\label{sec:p0-conclusion}

\begin{theorem}[Johnson 替换的 P0 状态]\label{thm:P0discharge}
假设~\ref{hyp:P0} 对显式 Johnson 表述成立。
\end{theorem}

\begin{proof}
保持 Heegaard 的 $93$-因子映射、七分量带 framing 的 AR link、六条与
$1513$ 条 Kirby 带、实际 $44$-通道切割、端点返回图与终端词比较提供所需桥梁、消去与 cabling 数据。
E13 链将所得闭图景与 $\Sigma_A^0$ 识别。
\end{proof}
\section{系数迹比较}\label{sec:c-comparison}

本节记录 Johnson 替换领圈的 C 论证。

\subsection{乘积 Hattori 等价}

紧凑词给出字母计数
\[
p_y=42+2=44,\qquad p_z=269+2=271.
\]
在循环 $m_2$ 词中，42 个 $y/Y$ 事件每一个都紧接其后跟一个不同的 $z/Z$ 事件；$r_{xy}$ 的两个 $y/Y$ 事件亦然。见证列出全部 44 对。
这些对在 P0 重构股沿其认证法向的 44 条乘积 ribbon 中实现，P0 立方体外另有 $227$ 个剩余 $z$-圆。范围是 Johnson 替换领圈，而非 $\partial W_2$ 中的嵌入切割。

令 $\mathcal C$ 为对象 $T:P_{86}\to P_{88}$ 的预分次有理 tangle 范畴，$F$ 为这些矩形构成的带框西-东 tangle。MWW 一柄定理将切割系数识别为
\[
M_R(T,T')=\KhR_2(R\cup T'\cup\overline T)
\]
及其所示移位与粘合作用
\cite[Theorem~4.7]{ManolescuWalkerWedrich2023}。枢轴 tangle 对偶与 $\Q$ 上 K\"unneth 给出此替换领圈上的分次同构
\begin{equation}\label{eq:HattoriActual}
H_{T,T'}:M_R(T,T')\xrightarrow{\cong}
\operatorname{Hom}_{\mathcal C}(FT,FT')\{-44\}
\otimes A^{\otimes227},
\qquad A=\Q[X]/(X^2).
\end{equation}
这些映射构成替换领圈上的系数双模同构。先在源 link 上粘合 $f:S\to T$ 再施加乘积同痕，与 rel 边界同痕于先施加同痕再以 $Ff$ 预复合。右作用是对后复合的同样论证。故两个作用方块都是同一同时曲面的粘合结合性。在链层面两个方块为
\[
\begin{array}{ccc}
C_R(T,T')&\xrightarrow{\ H_{T,T'}\ }&
C(FT,FT')\otimes C(S^1)^{\otimes227}\\
\big\downarrow {L_f}&&\big\downarrow {L_{Ff}\otimes1}\\
C_R(S,T')&\xrightarrow{\ H_{S,T'}\ }&
C(FS,FT')\otimes C(S^1)^{\otimes227},
\end{array}
\qquad
\begin{array}{ccc}
C_R(T,T')&\xrightarrow{\ H_{T,T'}\ }&
C(FT,FT')\otimes C(S^1)^{\otimes227}\\
\big\downarrow {R_g}&&\big\downarrow {R_{Fg}\otimes1}\\
C_R(T,U)&\xrightarrow{\ H_{T,U}\ }&
C(FT,FU)\otimes C(S^1)^{\otimes227}.
\end{array}
\tag{17}
\]
此处 $C_R$ 与 $C$ 为 strict Blanchet--Khovanov foam 复形。每个方块中两个复合都是在交换两个不相交领圈后得到同一粘合 foam。BHPW strict 函子性消去 projective 符号，取同调得 \eqref{eq:HattoriActual} 及 Johnson 替换领圈的两个 MWW 作用方块。方程~\eqref{eq:HattoriActual} 不是 $\partial W_2$ 中嵌入切割的链层面复形。

普通商陈述（含两个循环关系 span）在
形式化发展中核检查。\footnote{见
\texttt{Smooth4PC/RepresentableCoefficient.lean}。}

\begin{lemma}[实际系数双模]\label{lem:C1}
对 P0 替换，\eqref{eq:HattoriActual} 是实际 MWW 系数双模（含两个粘合作用）的自然同构。此外，下述映射 \(\operatorname{Sh}_q\) 是该双模的量子 trace，而非维数匹配的替代。
\end{lemma}

\begin{proof}
44 条乘积 ribbon 由 P0 重构股实现，具认证乘积法向 $z$-侧，给出 Johnson $y$-后-$z$ 配对
及 $227$ 个剩余 $z$-圆。\eqref{eq:HattoriActual} 的作用立方与第~\ref{sec:exact-algebra} 节的可表 $HH_0$ 约化完成比较。
\end{proof}

\subsection{量子 trace 与端点映射}

令 $R_q=\Q[q,q^{-1}]$。预分次循环关系为
\[
L_f(m)=q^{|f|}R_f(m).
\]
由引理~\ref{lem:C1}，\eqref{eq:HattoriActual} 与 counit 是 Johnson 替换领圈上的齐次系数态射。BPW 典范竖直-水平函子与 BHPW strict tangle/foam 函子给出
\begin{equation}\label{eq:actualShadow}
\operatorname{Sh}_q:q\operatorname{Tr}(\mathcal C;M_R)
\longrightarrow\operatorname{Hom}_{R_q}(E_{86},E_{88}),
\end{equation}
在 flat-base qHH$_0$/Chern 识别之后
\cite{BeliakovaPutyraWehrli2019,BeliakovaHogancampPutyraWehrli2023}。此处
\[
E_{86}=(V^{\otimes86})_{\lambda=86},\qquad
E_{88}=(V^{\otimes88})_{\lambda=86};
\]
前者秩为一，后者秩为 88。tangle 映射在限制到这些权空间之前是 $U_q(\mathfrak{sl}_2)$ 交织子。

基变换 $q\mapsto1+h$ 经有理局部化与 Noether 局部环完备分解，故为 flat。trace 余核对 $h$ 约化仅将 $q^{|f|}$ 换为 1，给出普通 MWW 系数 $HH_0$；张量积右正合性对此最后断言已足够。

\subsection{所选类与除法检测器}

令 $U:P_{86}\to P_{88}$ 为 P0 端点序中的所选单杯 tangle，并令 $T_1=F^{-1}U$。定义
\[
v_T=H_{T_1,T_1}^{-1}
  (\operatorname{Id}_U\otimes X^{\otimes227}).
\]
圆 counit 给出
$\operatorname{Sh}_h([v_T])=u_h$。其绝对次数为
\[
-44+227+315-4=494.
\]

在任何公开指标被指派之前，44 个实际乘积矩形给出
88 个物理端点：正侧记录实际 $y$ 源，负侧记录其配对的 $z$ 源。边界定向将两个
$r_{xy}$ wicket 按正向排序，将 42 个 $m_2$ wicket 按反向排序，每个矩形内先负侧后正侧。仅在此后公开位置表附着 Burau 指标。所得单项式映射 $P(q)$ 及其逆由这些物理标识符重算；全部枢轴系数均为
$+q^0$，且符号、次序与枢轴突变均失败。

相邻单缺陷基向量上归一化检查的 R-矩阵为
\[
\begin{pmatrix}1-q^{-2}&q^{-1}\\q^{-1}&0\end{pmatrix}.
\]
经位置基变换及 $t=q^{-2}$ 后，它是公开正 Burau 块；负块为其逆。物理缆具有混合定向。BPW 定向缆化公式将其在 writhe 依赖的分次移位与枢轴基映射下等同于全正模型。完整公开词 writhe 为零，故全局移位消失。任意次数零置换或枢轴坐标图 $P$ 同时传输算子、杯与盖：
\[
W\mapsto PWP^{-1},\qquad u\mapsto Pu,\qquad
\ell\mapsto\ell P^{-1}.
\]
矩阵系数不变；我们在该共同传输后选取公开端点标签。不使用未确定的 $\pm$ 常数项公式。在第~\ref{sec:exact-algebra} 节被运输的领圈坐标中，常数配对是那里记录的确定恒等式，而非符号 ansatz。

在 $\Q[[h]]$ 上定义
\[
D_h=\ell_h(\rho_h(W)-I)\operatorname{Sh}_h.
\]
这在量子系数 trace 上良定义，因为 \eqref{eq:actualShadow} 已施加量子循环性。精确 Magnus 计算与 Darn\'e 的 Andreadakis 定理给出 $W\in\Gamma_3(P_{44})$ \cite[Theorem~6.2]{Darne2021}。由此推出
\[
\rho_h(W)-I\in h^3\operatorname{End}(E_{88}).
\]
独立验证直接计算截断 Burau 表示的全部 $7{,}744$ 个矩阵条目；该计算不视为与 Andreadakis 成员关系等价，而是对同一三次消失的单独检查。
故 $D_h$ 取值于 $h^3\Q[[h]]$。除以 $h^3$ 唯一，对 $h$ 约化给出与提升无关的普通泛函
\[
D_3:HH_0(\mathcal C;M_R)\longrightarrow\Q.
\]
杯、盖与基变换中的正阶修正不改变三次系数。

四个混合指标 Burau 符号变体
\[
-53824,\quad -53760,\quad -59008,\quad -59072
\]
非零，但它们不是所选三次项。几何绑定的公开收据与
Lean 记录截断 Burau 标量
\begin{equation}\label{eq:selectedD3}
D_3([v_T])=2624
\end{equation}
在端点指标勘误之后。此整数目前仅是该
几何绑定的 Burau 值。它不闭合 C 或 S。

\subsection{完整二柄余锥}

\begin{lemma}[全缆常数项]\label{lem:C2}
对每个有限缆级及每个保号缆 braid \(b\)，端点算子的常数项等于其带符号置换。将该置换标准化后，剩余纯 braid 算子为 \(I+O(h)\)。
\end{lemma}

\begin{proof}
BPW 定向缆化作用由 strict BHPW tangle 函子得到。在 \(q=1\) 时，基本表示及其对偶上的 braiding 为普通对称交换，含相反定向因子的枢轴识别。故 braid 的常数项恰为其带符号端点置换的映射。保号纯 braid 置换为恒等，故常数项为 \(I\)。所有矩阵元为 \(q\) 的 Laurent 多项式，在 \(q=1\) 正则；在 \(q=1+h\) 后减去常数项得 \(h\) 的倍数。论证与缆副本数无关，故在每个有限 MWW 级成立。此陈述仅关于端点表示，不断言反平行 MWW beta 作用的未知对称群分解。
引理应用于引理~\ref{lem:C1} 之后的 Johnson 替换。
\end{proof}

仅断言 $D_3$ 在 C1 之后下降。本小节其余部分是引理~\ref{lem:C1} 之后的 C2 论证。

P0c 在重构股上提供乘积法向，从而在替换领圈中给出不相交平行级。我们首先在每一 MWW 缆态上类型化该构造。在 owner 序 $(m_2,m_3,r_{xy},r_{yz},r_{zx})$ 下，令
\[
n_y=(42,189,2,2,0),\qquad n_z=(269,1271,2,2,0).
\]
对 $r\in\mathbb N^5$，沿 $y$- 与 $z$-余核切割实际缆 $K(r,r)$ 得
\[
p_y(r)=\sum_i r_i n_{y,i},\qquad
p_z(r)=\sum_i r_i n_{z,i}.
\tag{24}
\]
由引理~\ref{lem:C1}，P0 矩形的平行副本给出与 (17) 同形的链同构 $H_r$，具相应端点集。$r_{zx}$ owner 与每个其它因子使用同一局部 counit 闭合；空自由词不被当作 split unknot。因所有副本位于同一固定管状乘积坐标图中，$H_r$ 与每个粘合作用及添加平行对可交换。

令 $s_0=e_{m_2}+e_{r_{xy}}$。若 $r\ge s_0$，令 $\Omega_r$ 为 $m_2$ 与 $r_{xy}$ 各选一个负物理副本与一个正物理副本的全部有限选择集。对 $\omega\in\Omega_r$，用 MWW 核心 counit 闭合每个未选乘积因子，用所选 44 个通道做点推，并记所得行的三次系数为 $D_{r,\omega}$。定义
\[
D_r=|\Omega_r|^{-1}\sum_{\omega\in\Omega_r}D_{r,\omega}.
\tag{25}
\]
对不在 $s_0$ 之上的态，用 once-dotted MWW 映射补缺失对，并通过拉回 (25) 定义 $D_r$。不相交支撑使这些添加的顺序无关紧要。故 $D_r$ 在 MWW 定义~3.1 的每个直和项上定义，含所选态之下的态。

对 beta，将 P0 领圈运动同时施加于所有物理副本。这是 \cite[Theorem~E]{BeliakovaPutyraWehrli2019} 的定向缆化作用，由 BHPW 严格化。常数置换重标 $\Omega_r$。放置标准化后，剩余纯 braid 由引理~\ref{lem:C2} 为 \(I+O(h)\)。算子、杯与盖的同时传输，连同 $\rho(W)-I=O(h^3)$，仅将行改变到四阶及以上。因此
\[
D_r\beta_i(b)=D_r
\tag{26}
\]
对每个 $b\in B_{r_i,r_i}$。此三次阶论证不使用反平行 braid 作用分解过对称群的断言，MWW 对此仍开放
\cite[Remark~3.3]{ManolescuWalkerWedrich2023}。

对一对添加，按新对是否被选中分裂目标放置。beta 移动交换所选新对与旧对，故 (26) 将两个子集约化到同一局部核心计算。在整个源上该计算为
\[
(\epsilon\otimes\epsilon)\Delta(1)=0,\qquad
(\epsilon\otimes\epsilon)\Delta(X)=1.
\]
对每个 owner（含 $r_{zx}$）使用同一局部 counit 方程；它们不由自由词消失推出。遗忘 distinguished 新对在放置集之间有常数大小纤维，(25) 中的归一化消去其基数。dotted 原始次数 $+2$ 被目标缆化移位 $-2$ 抵消。因此对每个 owner 与态，
\[
D_{r+e_i}\psi_i^{[0]}=0,\qquad
D_{r+e_i}\psi_i^{[1]}=D_r.
\tag{27}
\]
在 $s_0$ 之下第二个等式为定义；Frobenius 恒等式与不相交支撑交换证明第一个等式及所有 owner 方块可交换。

方程 (26)--(27) 连同已由 (17) 处理的系数 trace 关系，在整个缆化直和上证明
\[
D_3(\beta_i(b)x-x)=0,\quad
D_3(\psi_i^{[0]}x)=0,\quad
D_3(\psi_i^{[1]}x-x)=0.
\]
商 universal 性质与 MWW 定理~3.2 因此给出 $\overline D_3:\mathcal S^2_0(W_2;\Q)\to\Q$。

\subsection{C 结论与次数-$494$ 方块}

MWW 一柄与二柄映射嵌入交换图
\[
\begin{array}{ccccccc}
M_R(T_1,T_1)&\xrightarrow{H_{T_1,T_1}}&
\operatorname{End}(U)\{-44\}\otimes A^{\otimes227}
&\xrightarrow{\operatorname{id}\otimes\epsilon^{\otimes227}}&
\operatorname{End}(U)&\xrightarrow{d_3}&\Q\\
\big\downarrow{\mathrm{Glue}_{1h}}&&&&&&\big\Vert\\
\mathcal S^2_0(W_1;K(s_0,s_0))\{315\}
&\xrightarrow{\ \iota_{s_0}\{-4\}\ }&
\mathcal S^{2,0}_0(W_1;K)/(\beta,\psi)
&\xrightarrow[\cong]{\ \Phi\ }&\mathcal S^2_0(W_2)
&\xrightarrow{\overline D_3}&\Q.
\end{array}
\tag{28}
\]
左竖映射及其移位为 MWW 定理~4.7，底同构为 MWW 定理~3.2。上所选元为 $\operatorname{Id}_U\otimes X^{\otimes227}$。其逐次次数为
\[
-44+227=183,\qquad183+315=498,\qquad498-4=494.
\]
图 (28) 将其映到类 $x_2$，在 Johnson 替换领圈上满足 $\overline D_3(x_2)=2624\ne0$。

\begin{theorem}[系数比较 C]\label{thm:Cdischarge}
假设~\ref{hyp:P1} 对显式 Johnson 表述成立。
\end{theorem}

\begin{proof}
44 个实际 y/z 矩形与 227 个实际剩余圆运输给出引理~\ref{lem:C1} 中的系数双模识别；引理~\ref{lem:C2} 与 Reynolds 余锥验证两个作用及每个有限缆态。
\end{proof}

\paragraph{C/S 防火墙。}
C 中得到的 $B$-external naturality 涉及系数比较中的边界配边。它不把 $\partial W_2$ 中本质球面的局部作用等同于普通 dotted 球面求值；那是引理~\ref{lem:Sendpoint} 中的单独计算。
\section{三柄球面闭包}\label{sec:s-closure}

本节记录预期的 S 论证。闭流形 HJ
定理~5.3 不用于得出 $B$ 已固定。去掉最后的 $4$-柄，称
剩余柄体为 $W_3$。其边界为 $S^3$。反向三个
$3$-柄即在 $S^3$ 上附着三个三维 $1$-柄，故
\[
\partial W_2\cong\#^3(S^1\times S^2).
\]
实际 attaching 球面系的补连通，因为在该系上球面
surgery 给出 $W_3$ 的连通边界。

\begin{lemma}[更换完整球面系下的核不变性]
\label{lem:Ssystem}
令 \(\mathcal S\) 与 \(\mathcal A\) 为
\(\#^3(S^1\times S^2)\) 中的完整球面系。若它们经 isotopy、置换与
球面 slide 相关，则自 \(\mathcal S^2_0(W_2)\) 出发的两个总 MWW 三柄映射的核
一致。
\end{lemma}

\begin{proof}
attaching 球面的 isotopy 给出同构的柄附着。
置换仅重排不相交三柄。球面 slide 恰为
\(3\)--\(3\) 柄 slide。标准柄演算给出
结果三柄配边之间的微分同胚，在其入射边界 \(W_2\) 上为恒等。Skein-lasagna
映射的自然性因此给出 \(q_{\mathcal A}=\Phi q_{\mathcal S}\)，其中
\(\Phi\) 为目标模的同构。故它们的核一致。
\end{proof}

在 Johnson 替换反向
$1$-柄图景中选取标准完整系
\(\mathcal A=(A_1,A_2,A_3)\)，已与 P0 重构立方体 $B$ 不相交。
闭流形 HJ 定理~5.3 不用于得出 $B$ 已固定。
它仅与引理~\ref{lem:Ssystem} 一起使用：任何其他完整系，
含 going-up Kirby attaching 系，在闭流形中经 isotopy、
置换与 slide 相关，故核一致。
Horvat--Jab\l{}onowski 当前预印本 \cite{HorvatJablonowski2025}，
日期 2026 年 8 月 24 日，含引理~5.5（带点子球 slide 引理）与
引理~5.7（完整球面系的相对唯一性）。本文不调用这两条引理：核不变性仅使用其定理~5.3 与
引理~\ref{lem:Ssystem}。引理~5.7 针对具球面边界流形中的相对球面系，不是本节所用设定。
在反向 $1$-柄图景中，三个 belt 球面避开 P0 立方体，
dual 圈配对为恒等，且与 C1 剩余 link 及 C2 支撑不相交。该抽象 belt 模型的 $b=0$ foam 不是
在实际二柄核心圆盘恢复之前的 MWW 曲面。
记录的球面列给出代数下界计数
\[
(b_1,b_2,b_3)=(9920,1430,311).
\]
跟踪定向对偶圆盘词而不消去非相邻相反副本给出实际几何计数
\[
(b_1,b_2,b_3)=(12578,1824,409).
\]
其乘积法向副本轮廓与相应
$\epsilon^{\otimes b_j}\Delta^{b_j-1}$ 计算显式。
相对 PL 边界映射以层级表示：它施加全部 93 个
环境对偶圆盘因子，再施加两次 Kirby 边界映射的全部六条与 1513 条带，每个所得边界副本绑定到实际 owner。
矩阵本身不再是外部输入：将对偶 handlebody 的三个正方形压缩圆盘经全部 93 个环境
Johnson 因子运输，逐因子给出 $\wedge^2A=A^{-T}$，且
mapping torus 边界为 $I-A^{-T}$。单独的曲面运输
验证器重放两次 Kirby 领圈并拒绝缺失带。

MWW 对 $3$-柄商的内在重述使用
$\mathcal S^2_0(S^2\times D^2)$ 的局部模作用。在 $N=2$ 时，令 $A_0$ 为
带一个点的本质球面，$A_1$ 为无点球面。商
关系为
\[
A_0=1,\qquad A_1=0.
\]
若 $J_j$ 为 $A_j$ 的赤道，MWW 的 K\"unneth 识别给出
两半球映射的如下精确描述：
\[
\begin{array}{ccc}
\mathcal S^2_0(W_2;J_j)&\xrightarrow{\ (\Psi_{\Delta^+},\Psi_{\Delta^-})\ }&
\mathcal S^2_0(W_2)\oplus\mathcal S^2_0(W_2)\\
\big\downarrow\wr&&\big\downarrow{\ell\oplus\ell}\\
\mathcal S^2_0(W_2)\otimes\Q\{1,X\}
&\xrightarrow{\ (\ell\circ(1\otimes\epsilon),\,\ell\circ\mu_{A_j})\ }&
\Q\oplus\Q.
\end{array}
\tag{30}
\]
此处 $\epsilon(X)=1$，$\epsilon(1)=0$，而
$\mu_{A_j}(v\otimes X)=vA_0$ 且
$\mu_{A_j}(v\otimes1)=vA_1$。因此 S 等价于两个
全源恒等式
\[
\ell(vA_0)=\ell(v),\qquad \ell(vA_1)=0
\quad\text{对每个 }v\text{ 及 }j=1,2,3.
\tag{31}
\]
这将形式映射与 MWW
定理~3.7 中的 coequalizer 对等同，但不是实际本质球面的端点求值。

\begin{proposition}[固定检测器的球面替换]\label{thm:Sgeometry}
三柄总核可用 Johnson 替换反向 $1$-柄图景的标准 belt 球面系计算，该系避开
$B$，而 $B$ 中检测器保持固定。对此系，S 约化为
六个方程 (31)。
\end{proposition}

\begin{proof}
Johnson 替换的反向 $1$-柄图景放置三个
避开 $B$ 的 belt 球面，dual 圈的柄段与
owner belt 球面相交一次，且坐标图返回避开每个 belt 立方体。
闭流形 HJ 定理~5.3 将任何其他完整
系在闭流形中经 isotopy 与此系相关；引理~\ref{lem:Ssystem}
识别核。检测器留在那些 belt
球面避开的 $B$ 中。六个方程 (31) 于是为引理~\ref{lem:Sendpoint}。
\end{proof}

\subsection{本质球面端点比较}

\begin{lemma}[标准球面的端点分解]\label{lem:Sendpoint}
对每个 \(j\)、每个缆态及每个源元 \(v\)，
\(\mu_{A_j}(v\otimes a)\) 实际端点像的常数项是 \(v\) 的旧检测器像与
穿孔标准球面秩二 TQFT 映射的张量积。因此除法三次行满足
\[
\ell(vA_0)=\ell(v),\qquad \ell(vA_1)=0.
\]
\end{lemma}

\begin{proof}
使 \(A_j\) 与二柄余核横截，令 \(b_j\) 为
核心圆盘数。MWW 定理~3.10 的证明将
\(\Delta^-\) 半球（在核心圆盘恢复之前）表示为从 \(A_j\) 删去那些圆盘及
\(\Delta^+\) 圆盘后得到的连通亏格零曲面。在一柄切割后选取 generic movie。
因 \(A_j\) 与检测器片不相交，movie 的每个临界点完全属于两曲面之一。混合事件
因此仅为可逆 Reidemeister/枢轴移动，在端点由缆 braid 表示；无混合 birth、saddle 或 death。

在 \(q=1\) 时，引理~\ref{lem:C2} 将每个混合 braid 识别为
张量范畴的对称性。birth、saddle 与 death
对该对称性的自然性将所有混合 braid 移到 movie 末端。同一端点置换与枢轴坐标图传输
检测器行与曲面映射，故它们经同时
共轭抵消。故实际常数映射为
\[
\operatorname{Id}_{\rm old}\otimes\Delta^{b_j-1}\quad(b_j>0),
\]
当 \(b_j=0\) 时为 \(\operatorname{Id}_{\rm old}\otimes\epsilon\)。
恢复的二柄核心圆盘给出 \(b_j\) 个实际 counit。故
\[
\epsilon^{\otimes b_j}\Delta^{b_j-1}(1)=0,\qquad
\epsilon^{\otimes b_j}\Delta^{b_j-1}(X)=1
\tag{32}
\]
对 \(b_j>0\)，当 \(b_j=0\) 时解释为
\(\epsilon(1)=0,\epsilon(X)=1\)。
\(1,X\) 基由 MWW
例~3.8 中 \(\Delta^+\) cap 归一化，BHPW strict 函子性防止
\(\Delta^-\) movie 中的独立符号。

完整混合 braid 映射形如 \(P(I+O(h))\)。置换
\(P\) 刚被抵消，以 \(I+O(h)\) 预复合或后复合始于
$h^3$ 阶的行仅改变到 $h^4$ 阶及以上。
故常数计算 (32) 正是除法三次行所见的候选计算。对三个非零值
$b_j=12578,1824,409$，有限 Frobenius movie 分别有
$b_j-1$ 个余积鞍点与 $b_j$ 个 counit，并再现 (32)。故
几何 $\partial W_2$ 绑定存在；下述定理在全局接口打包完整 MWW 映射之前仍停留于所述断言边界。
\end{proof}

端点交换方块是上述引理~\ref{lem:Sendpoint} 的迭代余积--counit 模型，绑定到实际层级 $W_2$ 曲面。
\[
\begin{array}{ccc}
\mathcal S^2_0(W_2;J_j)&\longrightarrow&
\mathcal S^2_0(W_2)\oplus\mathcal S^2_0(W_2)\\
\downarrow&&\downarrow\\
\mathcal S^2_0(W_2)\otimes\Q\{1,X\}&\longrightarrow&\Q\oplus\Q,
\end{array}
\]
其下映射为实际 counit 行及引理~\ref{lem:Sendpoint} 中 movie 分解的实际 \(A_j\)-作用。

\begin{theorem}[相对三柄闭包 S]\label{thm:Sdischarge}
假设~\ref{hyp:P2} 对显式 Johnson 表述成立。
\end{theorem}

\begin{proof}
三个 H1 压缩圆盘经全部 93 个 Johnson 因子与全部 1519 条 Kirby 带运输。其几何边界词长为
$12578,1824,409$；所得穿孔球面 movie 在实际 $\partial W_2$ 上实现引理~\ref{lem:Sendpoint} 的映射。
\end{proof}
\section{最终识别}\label{sec:final-identifications}

MWW 定理~3.10 将迭代商与三柄后模等同，命题~3.4 给出四柄同构 \cite{ManolescuWalkerWedrich2023}。定理~\ref{thm:P0discharge}、定理~\ref{thm:Cdischarge} 与定理~\ref{thm:Sdischarge} 解除假设~\ref{hyp:P0}--\ref{hyp:P2}。对 Johnson 替换四柄图景 $X_J$，沿每个 belt 球面附着三柄恢复边界 $S^3$，附着 PL $4$-ball，标准球面模的 quantum degree-$494$ 分项在假设~\ref{hyp:P3} 证明中的计算连同 MWW 推论~3.5 而消失。实际 E13 柄链将 $X_J$ 与 $\Sigma_A^0$ 识别。

\begin{remark}[形式化边界]\label{rem:lean-boundary}
Lean 开发 \cite{Lean4} 将定理~\ref{thm:A} 证明为从结构 \texttt{ExternalGeometry} 与 \texttt{CSExternalGeometry} 出发的条件蕴含。它检验矩阵算术、几何绑定的 Burau 值 $2624$、次数 $494$ 与抽象商引理，但不构造几何实例。附录~\ref{app:lean-fields} 列出与 Lean 字段的对应。
\end{remark}
\input{sec-retired-assumptions-zh}
\section{相关工作}

Cappell--Shaneson 文献既有原始构造 \cite{CappellShaneson1976,AitchisonRubinstein1984}，也有大量标准性结果 \cite{Akbulut2010,Gompf2010,KimYamada2017,Iwaki2024}。我们对 Iwaki 的使用限于标准矩阵形、代表表与同伦球面判据。其表中 $(41,189,73)$ 条目是动机，而非奇点或异质性定理。

Manolescu--Neithalath 与 Manolescu--Walker--Wedrich 提供 skein lasagna 模的相关柄演算 \cite{ManolescuNeithalath2020,ManolescuWalkerWedrich2023}。BPW 量子 trace 函子性与严格 BHPW tangle 理论提供一般比较工具 \cite{BeliakovaPutyraWehrli2019,BeliakovaHogancampPutyraWehrli2023}。第 7 节将这些函子应用于实际替换系数并证明逐态自然性。

Ren 与 Willis 用 Khovanov 理论的 skein lasagna 方法区分其他异质四维流形 \cite{RenWillis2024}。其计算与 trace-73 构造无关，未用作其比较定理。

Horvat--Jab\l{}onowski 近期工作给出识别三柄附着系统的有用几何判据 \cite{HorvatJablonowski2025}。定理~5.3 仅与引理~\ref{lem:Ssystem} 用于闭流形中的核不变性，不用于固定 $B$。该预印本当前版本日期为 2026 年 8 月 24 日，含引理~5.5 与引理~5.7；二者均未调用。
\section{结论}

对 Johnson 替换，数学证书解除假设~\ref{hyp:P0}--\ref{hyp:P3} 并识别 $X_J\cong\Sigma_A^0$。Lean 开发仍只证明条件蕴含，因为分析性的 MWW 对象与映射尚未构造成 \texttt{ExternalGeometry} 的居民。故反例定理是有待专家审阅的纸面数学主张，尚非完全由 Lean 验证的定理。独立重放命令列于附录~\ref{sec:reproducibility-extended}。
\appendix

\section{已废止的紧凑与 Nielsen 路线}\label{app:retired-routes}

本附录记录在第~\ref{sec:johnson-presentation} 节采用 Johnson 替换之前曾检验并否决的路线。
此处材料均不用于正文论证。

符号 AR 公式与标准穿孔圆盘 braid 本身不能确定约化 handlebody 边界中嵌入的 $44$-passage 领圈。
显式 19 步 Nielsen 代表约化 $m_2$ 长度为 $309$、42 个 detector 通道，而紧凑代表长度为 $311$、44 个通道；因此它不是公开领圈词的来源。

对 $x^{-1}$ 的内共轭恢复 44 个带基点 passage 并通过 $F_3$ 的自由基检验，但不改变无基点循环脊类；它不是嵌入领圈见证。
对非内 IA 修正的有界搜索找到改变 detector 的 44 通道候选，但没有满足互补 handlebody 检验的符号置换 dual-meridian 扩张。

紧凑直脊词在 $F_3$ 上单射但不满射，故它们不能是保持 handlebody 的 $\psi_A\in\operatorname{Aut}(F_3)$ 的像。
正文中的 Johnson $\alpha_{ij}$ 分解规避了此阻碍。

在词层面，309 与 311 字母代表收集到同一正规形 $y^{40}z^{269}$，但所需的 framing 运输依赖于 $\operatorname{lk}(m_2,r_{yz})$，而词本身不能确定它。
为 Johnson CS 柄图构造的 railroad 约化 PD（引理~\ref{lem:P0d-link}）给出 $\operatorname{lk}(m_2,r_{yz})=0$。

\section{P0 重构协议}\label{app:p0-reconstruction}

可接受的 P0 证书由以下显式数据组成：三角化三球 $B$ 的有理 PL 数据、44 条高度单调带标号股 $c_i:[0,1]\to B$ 及其法向量、到约化 attaching link 的 passage 绑定、两个保持 framing 的局部消去 movie，以及有序 crossing movie。验证器检验 passage 绑定、不相交性、framing 运输，以及诱导相对 endpoint braid 与公开 $11340$ 字母词的逐字母相等。

\begin{lemma}[P0 证书的正确性]
设输入通过严格验证器，且其独立嵌入性、不相交性与局部 movie 收据对所声明 PL 复形有效。则 $B$ 是约化 AR handlebody 中嵌入的带框三球，其 $44$ 股构成带框领圈，相对 endpoint 几何 braid 等于公开 $11340$ 字母 pure braid。
\end{lemma}

\begin{proof}
PL 球与 passage 绑定映射将每条 $c_i$ 置于实际 AR 边界 chart 中，并在两次消去中保持 ownership 与 endpoint 顺序。高度单调性使 $c_i$ 成为领圈中的 braid；不相交收据使股两两不相交。法向量及其运输收据给出每股的 framing，并表明两次消去 movie 保持它。按高度递增读取认证 crossing 事件，得此嵌入 tangle 的相对 endpoint Artin 词。验证器最终相等将该词与公开词逐字母等同。因带标号 braid 的相对 endpoint 同痕类由其 Artin 词代表，所声称 braid 相等成立。
\end{proof}

\begin{remark}
无嵌入几何而导入 crossing 行的符号见证不满足本协议。第~\ref{sec:johnson-presentation} 节所用 Johnson 证书满足它。可机器检验的重放列于附录~\ref{sec:reproducibility-extended}。
\end{remark}
\input{sec-appendices-extra-zh}
\section{与 Lean 字段的对应关系}\label{app:lean-fields}

为检验形式化边界的读者，主要分配如下：
\begin{center}
\small
\begin{tabularx}{\textwidth}{@{}l X@{}}
\toprule
Lean 字段 & 数学义务 \\
\midrule
\texttt{x0}, \texttt{ell0}, \texttt{ell0\_x0} & P1 比较与所选 genuine MWW 类。 \\
\texttt{q01}, \texttt{ell1\_comp\_q01} & 完备二柄 beta/psi 商与下降。 \\
\texttt{q12}, \texttt{ell2\_comp\_q12} & S 的反向 belt 球面、MWW 半球映射与完备 coequalizer。非 E7「$B$ 已固定」断言。 \\
\texttt{transport}, \texttt{fourIso} & Johnson 替换 $X_J$ 的 E11，非 $\Sigma_A^0$。 \\
\texttt{s4DegreeZero} & 计算的 \texttt{EmptyKhQ}(494)=0 与来自 \texttt{certify\_t73\_e12\_s4.py} 的 PL $S^4$ Euler 数据。同构 $G(B^4)\simeq G(S^4)$ 与 $\mathcal S^2_0(S^4)\cong\mathbb Z$ 为 MWW 命题~3.4 与推论~3.5，引用而非形式化。仅空 link Lean 实例。 \\
\texttt{diffeomorphismEquiv} & lasagna 模的分次微分同胚不变性。 \\
\texttt{matrixConditionsToHomotopySphere} & E13 构造 CS 柄图；Lean \texttt{CSTopologyData} 不可居。 \\
\bottomrule
\end{tabularx}
\end{center}

第~\ref{sec:johnson-presentation}--\ref{sec:final-identifications} 节的几何输入已在论文中完成。它们尚未组装成 Lean 中 \texttt{ExternalGeometry} 的单一居民；这是形式化间隔，而非数学定理的假设。
\texttt{T73S4Inhabitant.lean} 中的空 link 控制仅打包 \texttt{EmptyKhQ}。Lean 定理不消费更强陈述，并将外部拓扑记录为结构字段而非内部形式化。

\bibliographystyle{amsplain}
\bibliography{references}
\end{document}
